\documentclass{article}
\usepackage{graphicx} % Required for inserting images

\usepackage[brazil]{babel}

\title{TCC}
\author{matheuscirillo }
\date{April 2026}

\begin{document}
    \maketitle

    \section{Introdução}
    \subsection{Equipes de Fórmula SAE}
    A Fórmula SAE é uma competição estudantil de âmbito mundial, organizada no Brasil
    pela SAE Brasil e realizada anualmente. As equipes são compostas por alunos
    de engenharia de uma mesma universidade, que ao longo do ano projetam,
    constroem e testam um protótipo de veículo monoposto de competição. Em 2019,
    a competição brasileira contou com mais de 70 equipes, evidenciando o crescimento
    e relevância da modalidade no cenário nacional. A competição é estruturada
    em provas estáticas e dinâmicas. As provas estáticas avaliam o projeto de engenharia,
    o plano de negócios e o detalhamento de custos e manufatura. As provas dinâmicas
    testam diretamente o desempenho do protótipo em situações como aceleração,
    skidpad, autocross e endurance — sendo esta última a de maior peso na pontuação,
    exigindo que o veículo percorra aproximadamente 22 km sem falhas. A equipe vencedora
    no Brasil garante ainda vaga para competir nos Estados Unidos no ano seguinte.
    Dado o calendário anual da competição, cada ciclo de desenvolvimento é único
    e irreversível: falhas durante as provas representam a perda de um ano
    inteiro de trabalho. Isso impõe à equipe uma exigência constante de robustez,
    confiabilidade e melhoria contínua em todos os subsistemas do veículo.
    \subsection{A Importância da Aquisição de Dados}
    Em ambientes de competição automotiva, entender o comportamento do veículo é
    condição essencial para extrair seu máximo desempenho. Sensores distribuídos
    ao longo do protótipo fornecem informações que permitem validar e aprimorar projetos
    de engenharia, orientar ajustes de setup, monitorar a integridade dos sistemas
    críticos em tempo real e comparar objetivamente o desempenho entre diferentes
    pilotos. Sem instrumentação adequada, a equipe depende exclusivamente do feedback
    subjetivo dos pilotos — informação valiosa, mas limitada e sujeita a
    variações de percepção. A aquisição de dados transforma essa subjetividade em
    grandezas mensuráveis: pressão de freio, curso de suspensão, ângulo de
    esterço, aceleração lateral, temperatura de óleo, entre outras. Isso permite
    identificar com precisão em qual trecho da pista o carro perde tempo, de que
    forma cada piloto utiliza os recursos do veículo e se as modificações realizadas
    entre testes produziram o efeito esperado. Em suma, o sistema de aquisição de
    dados é o principal elo entre o trabalho de engenharia realizado no galpão e
    as decisões tomadas na pista.
    \subsection{O Sistema de Aquisição da Equipe}
    O sistema de aquisição de dados da equipe foi consolidado ao longo de vários
    anos de desenvolvimento e testes. Sua arquitetura é baseada em dois módulos de
    aquisição com microcontroladores STM32, interligados por um barramento CAN a
    1 Mbit/s. Esse barramento conecta ainda os demais módulos eletrônicos do veículo
    — controle, transmissor e distribuição de energia — compondo uma rede coesa
    de nós de comunicação. Cada módulo de aquisição lê os sinais provenientes dos
    sensores analógicos e digitais instalados no carro — acelerômetro,
    giroscópio, sensores de pressão de freio, potenciômetros de suspensão e
    esterço, entre outros — e os transmite periodicamente pela rede CAN. Os
    dados são salvos localmente em cartão SD embarcado no próprio módulo e após
    cada sessão, os dados completos são extraídos do cartão SD e processados em um
    software dedicado, que gera arquivos estruturados para análise em
    ferramentas especializadas. Esse fluxo permite à equipe acompanhar a evolução
    do carro e dos pilotos ao longo do ano, embasando decisões de engenharia com
    informações objetivas e rastreáveis.
    % explain here what a formula sae team is
    % explain the importance of data validation
    % explain how data validation currently works for my team

    \section{O Problema}

    \subsection{Limitações do Sistema Atual}

    Embora o sistema de aquisição da equipe seja robusto e consolidado, sua arquitetura
    cabeada impõe restrições estruturais que limitam a capacidade de instrumentação
    do veículo. Essas limitações se manifestam em três cenários distintos, cada um
    com suas próprias implicações práticas.

    O primeiro cenário envolve componentes rotativos, como o interior das rodas.
    Nesses pontos, a conexão com o chicote elétrico do veículo é tecnicamente
    possível por meio de acopladores rotativos, porém a complexidade de instalação,
    o custo e a fragilidade desses componentes em ambiente de competição tornam
    a solução inviável na prática. Isso impede a leitura de sensores como temperatura
    de cubo ou deformação de manga de eixo — informações de alto valor para a dinâmica
    veicular.

    O segundo cenário diz respeito a regiões onde o cabeamento interfere diretamente
    no fenômeno a ser medido ou onde seu custo de instalação não se justifica pela
    frequência de uso. A asa traseira é o exemplo mais crítico: trata-se de um componente
    aerodinâmico cuja performance é sensível a qualquer perturbação no escoamento
    de ar ao redor de sua superfície. Além disso, testes de instrumentação na asa
    são pontuais e pouco frequentes, o que torna ainda menos justificável o
    investimento em um chicote dedicado — tanto em termos de material quanto de
    mão de obra para montagem e desmontagem a cada sessão.

    O terceiro cenário extrapola os limites físicos do veículo. Sensores posicionados
    fora do carro — como um sensor de passagem para detecção automática de voltas,
    ou equipamentos de medição estacionários utilizados em testes controlados —
    não têm como se comunicar com o sistema de aquisição embarcado por meios
    cabeados. Essa limitação impede a integração de dados externos ao fluxo de aquisição,
    obrigando a equipe a recorrer a registros manuais ou sistemas completamente
    separados, sem sincronização temporal com os demais canais do carro.
    % what are the limitations of our current setup
    % where they are more proeminent

    \section{Soluções Consideradas}

    Diante das limitações identificadas, três abordagens foram consideradas para
    expandir a capacidade de instrumentação do veículo.

    \subsection{Cabeamento Dedicado}

    A primeira opção consiste em aceitar os custos inerentes à arquitetura
    cabeada e investir na infraestrutura necessária para cada ponto de medição: acopladores
    rotativos para as rodas, chicotes dedicados para a asa traseira e extensões
    do chicote principal para outros pontos de difícil acesso. Essa abordagem preserva
    a arquitetura existente sem exigir nenhum desenvolvimento de software ou hardware
    adicional. No entanto, além dos custos de material e mão de obra já discutidos,
    ela não endereça o terceiro cenário levantado — sensores externos ao veículo
    — e continua impondo restrições físicas fundamentais que nenhum cabeamento é
    capaz de resolver.

    \subsection{Sistemas de Aquisição Independentes}

    A segunda opção consiste em desenvolver módulos de aquisição autônomos, cada
    um com sua própria memória local para gravação dos dados. Esses módulos operariam
    de forma completamente independente do sistema principal, sendo instalados
    nos pontos de interesse e recuperados após cada sessão para extração dos
    dados. Para viabilizar a análise conjunta, seria necessário algum mecanismo
    de sincronização temporal entre os módulos e o sistema embarcado no carro. Embora
    essa abordagem elimine a necessidade de cabeamento, ela introduz uma nova
    complexidade operacional: a sincronização entre sistemas independentes é propensa
    a erros, e os dados coletados não estariam disponíveis em tempo real, o que limita
    seu uso durante a sessão de testes.

    \subsection{Solução Sem Fio}

    A terceira opção — e a adotada neste projeto — é o desenvolvimento de uma
    ponte de comunicação sem fio entre os pontos de medição e o sistema de
    aquisição existente. Nessa abordagem, um nó sensor captura os dados
    localmente e os transmite por rádio para um nó receptor conectado ao
    barramento CAN do veículo. O receptor injeta os dados na rede como se fossem
    provenientes de qualquer outro nó cabeado, tornando a comunicação sem fio
    completamente transparente para o restante do sistema.

    Essa solução elimina as restrições de cabeamento, preserva o fluxo de dados existente
    sem alterações na arquitetura do sistema principal e abre espaço para instrumentação
    em pontos até então inacessíveis — incluindo componentes rotativos, regiões
    aerodinâmicas sensíveis e sensores externos ao veículo. A principal restrição
    a ser endereçada no desenvolvimento é a garantia de confiabilidade na entrega
    dos dados, uma vez que protocolos de comunicação sem fio, por natureza, não
    oferecem as mesmas garantias de integridade que uma conexão física.
    % the possible solutions we arrived
    % introduce the ideia of a wireless solution

    \section{Requisitos}

    Para que a solução sem fio seja viável em ambiente de competição, ela deve
    satisfazer um conjunto de requisitos funcionais e não funcionais que
    refletem tanto as restrições operacionais da equipe quanto as exigências
    técnicas de integração com o sistema existente.

    \subsection{Desempenho e Latência}

    O sistema deve ser capaz de transmitir dados com latência compatível com as frequências
    de aquisição já praticadas pela equipe. Atrasos na entrega dos dados
    comprometem a sincronização temporal com os demais canais do sistema e podem
    tornar a informação inutilizável para análise.

    \subsection{Conexão Instantânea}

    O sistema não pode depender de um processo de estabelecimento de conexão entre
    os nós antes do início da transmissão. Em cenários como um sensor de
    passagem externo ao veículo, o carro pode estar em alta velocidade quando passa
    pelo ponto de medição — qualquer overhead de handshake inviabilizaria a coleta
    do dado. A comunicação deve estar sempre pronta para transmitir, sem negociação
    prévia entre os dispositivos.

    \subsection{Múltiplas Conexões Simultâneas}

    O sistema deve suportar múltiplos nós sensores transmitindo simultaneamente
    para um ou mais receptores. Isso é especialmente relevante para a instrumentação
    das quatro rodas do veículo, onde todos os sensores precisam operar em
    paralelo e ter seus dados recebidos e injetados na rede CAN de forma
    coordenada.

    \subsection{Transparência com o Sistema CAN Existente}

    A integração com o barramento CAN deve ser feita com o mínimo de modificações
    possível na arquitetura existente. O sistema sem fio deve preservar os
    identificadores CAN originais de cada mensagem e manter as verificações de
    integridade nativas do protocolo — incluindo CRC e acknowledgment —
    garantindo que o restante da rede trate os dados recebidos sem fio de forma
    idêntica aos dados provenientes de nós cabeados.

    \subsection{Requisitos Físicos e de Custo}

    Os nós sensores devem ser fisicamente compactos e leves, uma vez que serão
    instalados em regiões onde massa e volume são restrições relevantes — como
    rodas e componentes aerodinâmicos. O custo total do sistema deve ser compatível
    com o orçamento de uma equipe universitária, priorizando componentes
    acessíveis sem comprometer os requisitos funcionais.

    \subsection{Modularidade}

    A arquitetura do sistema deve ser modular, permitindo adicionar novos nós
    sensores sem alterações na infraestrutura existente. Cada novo ponto de medição
    deve poder ser incorporado ao sistema de forma independente, com impacto
    mínimo nos demais nós e sem necessidade de reconfiguração do receptor ou do
    barramento CAN.
    % what are the requirements for a successful system

    \section{Framework}
    \subsection{Microcontrolador}

    A primeira decisão de hardware foi a escolha do microcontrolador para o nó sensor.
    A opção mais natural inicialmente seria utilizar a família STM32, já presente
    no sistema de aquisição da equipe e fornecida por um patrocinador. % TODO: confirmar patrocínio

    A linha STM32 não possui nenhum chip com Wi-Fi nativo. A única família com comunicação
    sem fio integrada é a STM32WB, que oferece Bluetooth Low Energy e 802.15.4 —
    sem suporte a Wi-Fi. O STM32WB integra um núcleo Cortex-M4 para a aplicação
    e um Cortex-M0+ dedicado ao stack de rádio, sendo uma plataforma robusta, bem
    documentada e familiar à equipe — uma opção forte, desde que o BLE atendesse
    aos requisitos de comunicação.

    Para contornar a ausência de Wi-Fi nativo, uma alternativa seria combinar um
    STM32 convencional com um módulo Wi-Fi externo. No entanto, essa abordagem exigiria
    três componentes distintos na placa do nó sensor: o microcontrolador, o módulo
    de comunicação e um circuito de gerenciamento de bateria separado. A
    combinação resulta em uma solução com footprint físico, complexidade de integração
    e custo de desenvolvimento incompatíveis com os requisitos de compacidade e leveza
    do projeto, e foi descartada por esse motivo.

    A terceira opção considerada foi o ESP32, plataforma da Espressif que integra
    nativamente Wi-Fi e Bluetooth em um único chip. Módulos de desenvolvimento
    compactos e com gerenciamento de bateria integrado estão amplamente
    disponíveis para essa plataforma — como o Beetle ESP32-C6 da DFRobot, de apenas
    25×20,5 mm e 17 g, que inclui o circuito de carga Li-ion na própria placa, sem
    necessidade de componentes externos adicionais.

    A decisão entre as duas plataformas foi tomada em conjunto com a escolha do protocolo
    de comunicação, discutida a seguir.

    \subsection{Protocolo de Comunicação Sem Fio}

    As opções naturais de protocolo foram Bluetooth Low Energy, Wi-Fi e, posteriormente,
    ESP-NOW.

    O BLE, disponível tanto no STM32WB quanto no ESP32, foi a primeira opção
    considerada. No entanto, o protocolo apresenta limitações fundamentais
    frente aos requisitos do sistema. O estabelecimento de uma conexão BLE envolve
    um processo de descoberta e pareamento que introduz latência inicial incompatível
    com o requisito de conexão instantânea — um sensor externo ao veículo, por
    exemplo, não teria tempo de completar esse handshake antes de o carro passar
    pelo ponto de medição. Além disso, o gerenciamento de múltiplas conexões BLE
    simultâneas adiciona complexidade considerável à aplicação. Essas limitações
    descartaram o BLE como protocolo viável e, consequentemente, tornaram o
    STM32WB uma opção menos atrativa, uma vez que seu diferencial principal não poderia
    ser aproveitado.

    O Wi-Fi resolve o problema de múltiplas conexões, mas mantém o overhead de
    associação à rede, o que continua incompatível com o requisito de transmissão
    instantânea.

    Durante o desenvolvimento, foi identificado o ESP-NOW, um protocolo proprietário
    da Espressif projetado especificamente para comunicação direta entre dispositivos
    ESP32, sem infraestrutura de rede e sem necessidade de estabelecimento de
    conexão. O ESP-NOW opera sobre a camada física do Wi-Fi, herdando sua
    verificação de integridade por CRC em nível de frame, e permite que um dispositivo
    transmita dados para outro imediatamente, sem qualquer negociação prévia. Essa
    característica, combinada com o suporte nativo a transmissões em broadcast
    para múltiplos receptores, atende diretamente aos requisitos de conexão instantânea
    e múltiplos nós.

    O protocolo é deliberadamente minimalista: não oferece confirmação de
    entrega nativa em modo broadcast, o que significa que a confiabilidade de entrega
    precisa ser construída na camada de aplicação. Esse trabalho adicional foi
    considerado um custo aceitável, dado que o ESP-NOW resolve de forma elegante
    os demais requisitos do sistema — e sua disponibilidade exclusiva na plataforma
    ESP32 consolidou a escolha do ESP32 como microcontrolador do projeto.
    % hardware selection, what were the options and why we choose what we choose
    % wireless protocol selection

    \section{Desenvolvimento}
    A integração de uma ESP32 ao barramento CAN do veículo — leitura dos sinais do
    transceiver, inicialização do driver TWAI e transmissão/recepção de frames —
    é um procedimento amplamente documentado e foi implementada sem dificuldades
    relevantes. O módulo \texttt{can.cpp} encapsula essa camada: inicializa o driver
    com os pinos configuráveis, gerencia estados de erro e recuperação do
    barramento, e expõe funções simples de envio e recepção. Por se tratar de
    uma etapa bem estabelecida, o restante desta seção concentra-se no desenvolvimento
    da camada de aplicação sem fio — o protocolo WCAN — que constitui a contribuição
    principal deste trabalho.

    \subsection{Arquitetura Geral}

    O sistema WCAN foi estruturado em quatro módulos de software com
    responsabilidades bem definidas: \texttt{wcan.cpp} (inicialização e callbacks
    ESP-NOW), \texttt{wcan\_sender.cpp} (fila de envio e mecanismo de retransmissão),
    \texttt{wcan\_receiver.cpp} (fila de recepção, filtragem e confirmação) e \texttt{wcan\_utils.cpp}
    (codificação/decodificação de pacotes e gerenciamento de peers). A separação
    em módulos permite que cada aspecto do protocolo seja desenvolvido e testado
    de forma independente.

    A unidade fundamental de dados do sistema é o \texttt{data\_packet\_t}, que encapsula
    o identificador CAN da mensagem original (\texttt{can\_id}), um contador de ticks
    do FreeRTOS (\texttt{tick\_count}) utilizado para correlacionar pacotes com
    seus respectivos acknowledgments, o endereço MAC do remetente e o payload de
    dados em tamanho variável. Esse pacote é serializado pelo módulo de
    utilidades em um buffer de bytes para transmissão via ESP-NOW e
    desserializado no lado receptor, mantendo a semântica CAN intacta ao longo de
    todo o caminho sem fio.

    \subsection{Desacoplamento por Filas}

    Uma decisão de projeto central foi o desacoplamento entre os callbacks do
    ESP-NOW e o processamento efetivo dos dados. O callback de recepção do ESP-NOW
    é executado no contexto de interrupção do Wi-Fi, um ambiente com restrições
    severas de tempo de execução e de operações permitidas. Realizar qualquer processamento
    significativo nesse contexto — como filtragem, reenvio de acknowledgments ou
    invocação de callbacks do usuário — comprometeria a responsividade do rádio
    e poderia causar perda de pacotes subsequentes.

    Para resolver essa restrição, o sistema utiliza duas filas do FreeRTOS: uma
    fila de envio (\texttt{send\_queue}) e uma fila de recepção (\texttt{recv\_queue}).
    O callback de recepção do ESP-NOW limita-se a desserializar o pacote recebido,
    classificá-lo como dado ou acknowledgment, e encaminhá-lo para a fila apropriada.
    Duas tasks dedicadas — \texttt{SendProcessingTask} e \texttt{RecvProcessingTask}
    — consomem suas respectivas filas em contexto de task normal, onde não há restrições
    de tempo ou de chamadas de sistema.

    Essa arquitetura produtor-consumidor garante que o callback de interrupção execute
    em tempo mínimo e constante, independentemente da complexidade do
    processamento necessário, e permite que o sistema escale para múltiplos sensores
    transmitindo simultaneamente sem degradação do tratamento de interrupções.

    \subsection{Mecanismo de Confiabilidade}

    Conforme discutido na seção anterior, o ESP-NOW em modo broadcast não oferece
    confirmação de entrega nativa. O sistema WCAN implementa um mecanismo
    próprio de acknowledgment e retransmissão para suprir essa lacuna.

    Quando um nó sensor envia um pacote de dados, a \texttt{SendProcessingTask}
    transmite o pacote via broadcast e inicia um timer periódico de
    retransmissão. Um semáforo binário (\texttt{send\_semaphore}) bloqueia o
    processamento do próximo pacote da fila até que uma das duas condições seja
    satisfeita: o recebimento de um acknowledgment correspondente ou o esgotamento
    do número máximo de tentativas de reenvio. Essa serialização garante que
    apenas um pacote esteja em trânsito por vez, simplificando o rastreamento de
    acknowledgments.

    No lado receptor, ao receber um pacote de dados válido, a \texttt{RecvProcessingTask}
    envia imediatamente um acknowledgment unicast de volta ao remetente, contendo
    o \texttt{tick\_count} do pacote original. O remetente utiliza esse \texttt{tick\_count}
    para correlacionar o acknowledgment com o pacote pendente — uma vez que
    múltiplos receptores podem enviar acknowledgments para o mesmo pacote
    broadcast. Ao receber o acknowledgment correspondente, o semáforo é liberado
    e o próximo pacote da fila é processado.

    Caso nenhum acknowledgment seja recebido dentro do intervalo configurado (\texttt{WCAN\_RETRY\_DELAY}),
    o timer dispara automaticamente o reenvio do pacote, repetindo o processo até
    o limite de tentativas (\texttt{WCAN\_MAX\_RETRY\_COUNT}). Se o limite for atingido
    sem confirmação, o pacote é descartado e o sistema avança para o próximo, registrando
    o erro via log.

    \subsection{Filtragem de Mensagens}

    Para atender ao requisito de comunicação N:M — onde múltiplos sensores
    transmitem em broadcast e múltiplos receptores escutam — o sistema
    implementa um mecanismo de filtragem por identificador CAN no lado receptor.
    Na inicialização, cada nó receptor pode opcionalmente fornecer uma lista de
    CAN IDs permitidos. A função \texttt{FilterData}, chamada para cada pacote recebido
    que não seja um acknowledgment, verifica se o \texttt{can\_id} do pacote consta
    na lista de IDs permitidos. Pacotes com identificadores fora da lista são
    descartados silenciosamente, sem gerar acknowledgment nem ocupar espaço na fila
    de recepção. Quando a filtragem está desativada, todos os pacotes são aceitos.

    Esse mecanismo permite que cada receptor processe apenas os dados relevantes
    para sua função, mesmo em um ambiente de broadcast onde todos os pacotes são
    visíveis a todos os nós — comportamento análogo ao filtro de aceitação
    presente em controladores CAN por hardware.

    \subsection{Integração com o Barramento CAN}

    O módulo \texttt{main.cpp} demonstra a integração completa do sistema. A aplicação
    identifica o papel de cada ESP32 pelo seu endereço MAC: nós sensores
    inicializam o WCAN sem filtragem e criam tasks de leitura de dados que
    alimentam a fila de envio; nós receptores inicializam tanto o WCAN (com
    filtragem habilitada) quanto o driver CAN, e definem um \texttt{RecvCallback}
    que, ao receber um pacote via WCAN, retransmite seu conteúdo no barramento CAN
    local utilizando o identificador CAN original. Do ponto de vista dos demais
    módulos do veículo conectados ao barramento, os dados provenientes dos sensores
    sem fio são indistinguíveis de dados originados em nós cabeados — preservando
    a compatibilidade total com o sistema de aquisição existente.

    \section{Test results}
    % ainda preciso definir melhor quais testes rodar pra colocar aqui, aceito sugestões
    % testes de rampa visandogarantir que a rampa chega inteira do outro lado, roda umas 50x e ve o quanto foi perdido
    % teste 1:1
    % teste N:1
    % teste N:M
    % eu tenho 5 ESP's em mão
\end{document}